%==============================================================================
% design-fontes.tex — Tipografia XeLaTeX (fontspec)
% Corpo de leitura moderno (Noto Serif). Pré/pós-textuais trocam para
% Liberation Serif via hooks em design-margem.tex.
%==============================================================================
\setmainfont{Noto Serif}[Ligatures=TeX]
\setsansfont{Inter}[Ligatures=TeX]
\setmonofont{DejaVu Sans Mono}[Scale=0.92]

% Famílias auxiliares
\newfontfamily\abntserif{Liberation Serif}[Ligatures=TeX] % ABNT pré/pós-textual
\newfontfamily\tituloface{Inter}[Ligatures=TeX]           % títulos/aberturas

%------------------------------------------------------------------
% Letra capitular (drop cap) colorida pelo nível do capítulo
% Uso: \capitular{nível 0-5}{primeira letra}{resto da 1ª palavra}
% A cor da capitular sinaliza o nível (verde=0 ... roxo=PhD).
%------------------------------------------------------------------
\newcommand{\capitular}[3]{%
  \lettrine[lines=3, findent=2pt, nindent=0pt]%
    {\tituloface\bfseries\textcolor{\cornivel{#1}}{#2}}{#3}}
